\section{The sieve lower bound}\label{sec:density}

The quantitative route to Conjecture~\ref{goldbach:conj} replaces the
existence question by a counting one.  By Lemma~\ref{reflect:lem} the
\term{goldbach-count} is the size of a set intersection, so any lower bound on
that size that is eventually positive settles the conjecture for all large $n$,
leaving a finite range covered by Theorem~\ref{verify:thm}.

\begin{definition}\label{sifted:def}
For even $n$ and a parameter $z\geq2$, the \dfnas{sifted-count}{sifted count}
$\gc_{z}(n)$ is the number of $p\leq n$ such that neither $p$ nor $n-p$ has a
prime factor below $z$.
\end{definition}

Sieve methods bound $\gc_{z}(n)$ from above and below for $z$ a small power of
$n$, and $\gc(n)\leq\gc_{z}(n)$ for every $z$, which is the wrong direction.
The route is therefore to control the difference.

\begin{prob}\label{density:prob}
Show that there is a constant $c>0$ with
\[
  \gc(n)\ \geq\ \frac{c\,n}{(\log n)^{2}}
\]
for every sufficiently large even $n$.
\end{prob}

A positive answer implies Conjecture~\ref{goldbach:conj} for all large $n$,
and hence, with Theorem~\ref{verify:thm}, in every range that has been checked
below the implied threshold.

\begin{prop}\label{upper:prop}
The bound of Problem~\ref{density:prob} is of the correct order: there is a
constant $C$ with $\gc(n)\leq C\,n/(\log n)^{2}$ for every even $n\geq4$.
\end{prop}

\begin{proof}[Proof sketch]
This is the standard upper-bound sieve applied to the sequence
$\{p(n-p)\}$; it is classical, and is recorded here only to fix the shape of
the target.  A complete proof is not reproduced, so this statement stays in the
companion rather than the results edition.
\end{proof}

\emph{Why the parity obstruction bites.}  A pure sieve cannot distinguish
numbers with an even number of prime factors from those with an odd number, and
the \term{sifted-count} counts both.  Any argument that resolves
Problem~\ref{density:prob} must therefore import information beyond the sieve
axioms; the \termas{hardy-littlewood}{Hardy--Littlewood} heuristic predicts the
constant but supplies no such input.  Numerical support is in
\texttt{tools/python/goldbach\_counts.py}.
